\documentclass[10pt,a4paper]{amsart}
\usepackage[margin=19mm]{geometry}
\usepackage{amsmath,amssymb,mathtools}
\usepackage[scr=rsfs,cal=euler]{mathalfa}
\usepackage{xfrac}
\usepackage[unicode,hidelinks,bookmarks=false]{hyperref}
\newcommand{\ie}{i.e.\ }
\newcommand{\cf}{cf.\ }
\newcommand{\resp}{respectively\ }
\newcommand{\Subsection}{\subsection}
\providecommand{\nobreakdash}{\nobreak\hbox{-}}
\setlength{\emergencystretch}{2em}
\multlinegap0pt
\usepackage{xcolor}
\usepackage{amssymb}
\usepackage{tikz}
\newtheorem{prop}{Proposition}
\title{All $4\times4$ solutions of the quantum Yang--Baxter equation}
\author{Marius de Leeuw, Vera Posch}
\date{Source: arXiv:2411.18685v6 (CC BY 4.0)}
\begin{document}
\maketitle

\noindent\emph{Source and licence.} All $4\times4$ solutions of the quantum Yang--Baxter equation. Version arXiv:2411.18685v6 (CC BY 4.0), distributed by arXiv under the Creative Commons Attribution 4.0 International licence, http://creativecommons.org/licenses/by/4.0/, as recorded in the arXiv licence metadata for this exact version and shown on the abstract page https://arxiv.org/abs/2411.18685v6. Full licence notice: \texttt{licenses/arxiv-2411.18685-CC-BY-4.0.txt}.\par\smallskip

\noindent\textit{Editorial scope.} Counted unit: the article's prop lem:RegularRLL (source0.tex lines 642--649). The article's complete original proof of it is delivered with it (source0.tex lines 650--695, one \texttt{proof} environment of 46 lines). No mathematics is written, repaired or re-derived: the statement and the proof are the article's own lines, in the article's own notation, and no retained line is substituted or re-flowed.  No further source-local context is needed: every cross-reference of every retained line resolves inside the two delivered spans.  Nothing was omitted from any retained span, so the package carries no omission record.  No retained line contains a citation, so the wrapper carries no bibliography and no bibliography entry.  No retained line contains a figure environment, an image-inclusion command or drawing code, so no image is delivered and the package declares no asset dependency.  Editorial formatting only, in addition: an AMS \texttt{amsart} preamble that restates the article's own declarations and macros used by the retained lines, namely the environments \texttt{prop} on separate counters, together with the standard packages those lines need.  The retained environments print in this document as Proposition 1 in order of appearance, whereas the article prints the same items with the numbering of its own full document.  The complete native source of the article as distributed in the arXiv e-print for this exact version is bundled unchanged as \texttt{sources/169250/source0.tex}.
\begin{prop}\label{lem:RegularRLL}
Let $L$ be a regular Lax operator which satisfies the fundamental commutation relations. Then the corresponding $R$ matrix is regular as well
%
\begin{align}
R(u,u) = P.
\end{align}
% 
\end{prop}

\begin{proof}
By evaluating the RLL relations at $u=v=0$, we find that
%
\begin{align}
R(0,0) \sim P.
\end{align}
%
Since $R$ is fixed up to an overall normalisation, we can set $R(0,0)=P$ without loss of generality. Now let us expand the Lax operator  and the $R$-matrix as
%
\begin{align}\label{eq:Lexpand}
&L(u) = P(1+ u \mathcal{H} + \sum_{i=2}^\infty u^i \check{L}^{(i)} ),
&& R(u,u) = P(1+ \sum^\infty_{i=1} u^i \check{R}^{(i)} ).
\end{align}
%
We plug this into the RLL relations and set $v=u$ to find
%
\begin{align}
&P_{ab}(1+ u \check{R}_{ab}^{(1)} + \ldots) P_{an}(1+u \mathcal{H}_{an} +\ldots) P_{bn}(1+u \mathcal{H}_{bn} +\ldots) =\nonumber  \\
 &\qquad \qquad P_{bn}(1+u \mathcal{H}_{bn} +\ldots) P_{an}(1+u \mathcal{H}_{an} +\ldots) P_{ab}(1+ u \check{R}_{ab}^{(1)} + \ldots).
\end{align}
%
At leading order this is satisfied by construction and at the next order we find
%
\begin{align}
&P_{ab}  \check{R}_{ab}^{(1)} P_{an} P_{bn} + P_{ab}  P_{an} \mathcal{H}_{an}P_{bn} + P_{ab} P_{an} P_{bn} \mathcal{H}_{bn}
=\nonumber  \\
 &\qquad \qquad 
 P_{bn}P_{an}  P_{ab}  \check{R}_{ab}^{(1) }+  P_{bn}  P_{an} \mathcal{H}_{an}P_{ab} +  P_{bn} \mathcal{H}_{bn}P_{an}P_{ab} ,
\end{align}
%
from which it follows that $\check{R}_{ab}^{(1)} =\check{R}_{bn}^{(1)}$ which can only be true for $\check{R}_{ab}^{(1)} \sim 1$. Hence this coefficient can be put to 0 by choosing again an appropriate normalisation of $R$. One can then recursively show that all orders in $R$ can be chosen to vanish. More specifically, by cancelling the permutation operators in the RLL relations we find
%
\begin{align}
&(1+ \sum^\infty_{i=1} u^i \check{R}^{(i)}_{bn} )(1+u \mathcal{H}_{ab} +\ldots)(1+u \mathcal{H}_{bn} +\ldots) 
=\nonumber\\
&\qquad\qquad (1+u \mathcal{H}_{ab} +\ldots) (1+u \mathcal{H}_{bn} +\ldots) (1+ \sum^\infty_{i=1} u^i \check{R}^{(i)}_{ab} ).
\end{align}
%
This yields an equation for $\check{R}^{(i)}$ in terms of the lower order expansion terms of the form
%
\begin{align}
\check{R}^{(i)}_{bn} + \sum_{k+l+m=i} \check{R}^{(k)}_{bn}\check{L}^{(l)}_{ab}\check{L}^{(m)}_{bn} = \check{R}^{(i)}_{ab} + \sum_{k+l+m=i} \check{L}^{(l)}_{ab}\check{L}^{(m)}_{bn} \check{R}^{(k)}_{ab},
\end{align}
%
where the indices in the sum run over positive integers, \textit{i.e.} $l,m>0$. By the induction hypothesis we have that $\check{R}^{(k)}=0$ for $k<i$, from which it follows that $\check{R}^{(i)}_{bn} = \check{R}^{(i)}_{ab}$. This again implies that $\check{R}^{(i)}\sim1$ and can be set to 0 by an appropriate normalisation.
\end{proof}

\end{document}
